% ===========================================================================
% Chapter 8: Verification
% ===========================================================================
\chapter{Verification}
\label{ch:verify}

No number in this report was typed in by hand. This chapter shows how numbers are
produced and checked, how to read a registry key, and what the checks catch.

% ===========================================================================
\section{The registry and its keys}
\label{sec:keys}

The script \texttt{verify/recompute\_25m0328.py} reads the model output committed to the
repository, recomputes every quantity the project uses, and writes them to
\texttt{verify/registry.json}: 10830 named values in 71 groups. Anyone can rebuild it
from the committed output by running the script, which also states the conventions of
Chapter~\ref{ch:conv}.

Each value has a \emph{key}: parts separated by full stops, read from left to right,
from the general to the particular,
\emph{group}\,.\,\emph{case or species}\,.\,\emph{quantity}\,.\,\emph{qualifier}. The
group names the source, usually an experiment; the middle parts narrow it to one run,
species, reaction or window; the last part names the quantity and, where needed, a
wavelength or a unit. Not every key has all four parts. Table~\ref{tab:keygroups} lists
the groups quoted in this report, and Appendix~\ref{app:filemap} lists, for each
experiment, its run files, scripts and registry groups.

\begin{table}[htbp]
  \centering
  \caption[Registry groups quoted in this report]{Registry groups quoted in this report.}
  \label{tab:keygroups}
  \footnotesize
  \setstretch{1.0}
  \begin{tabularx}{\textwidth}{>{\raggedright\arraybackslash}p{3.3cm}
                                >{\raggedright\arraybackslash}X}
    \toprule
    Group & What it holds \\
    \midrule
    \texttt{e1} to \texttt{e14} & One experiment each; letters mark a case or sub-study
      (\texttt{e7a} is Experiment~7A, \texttt{e3fp} the whole-field test of Experiment~3 at all four amplitudes) \\
    \texttt{res} & Largest change in a sulfur reservoir species, per case \\
    \texttt{dJ}, \texttt{dJy}; \texttt{dJz}, \texttt{dJyz} & Photolysis change as
      $\Delta J$ and as $\Delta(Jn)$; the heights where each peaks \\
    \texttt{nom}, \texttt{conv} & Reference and baseline runs; steady-state tests \\
    \texttt{pnom}, \texttt{p9}, \texttt{spec\_*}, \texttt{tau} & Absorber profiles,
      spectra, and optical depths built from them \\
    \texttt{scr}, \texttt{iso}, \texttt{corr} & Experiment~8 screen, Experiment~11
      isomers, spectral correlations \\
    \texttt{five}, \texttt{share320400} & Comparison of the five spectra; share of
      absorption between 320 and 400~nm \\
    \texttt{falselead} & The optical depths of the false lead (Chapter~3) \\
    \texttt{id}, \texttt{chr}, \texttt{chr8}, \texttt{twoby2}; \texttt{const},
      \texttt{p2} & Tests of
      Experiments~9 and~10; fixed constants; Phase~2 mixing times \\
    \bottomrule
  \end{tabularx}
\end{table}

\textbf{Example 1.} \texttt{e12.S2O2.short.360}: Experiment~12; species S\sub{2}O\sub{2}
(OSSO); shortfall factor; at 360~nm. Its value, $2.4563\times10^{5}$: the model holds that many
times too little OSSO above 69~km.

\textbf{Example 2.} \texttt{res.e7a.SO.inband}: reservoir responses; case 7A, the absorber
moved up; species SO; largest relative change inside the response band. Its value,
8.404\%, is one of the two abundance changes above the chemistry floor.

\textbf{Example 3.} \texttt{e7.7a.SO2.R781.col} and its twin ending \texttt{col\_km}:
Experiment~7; case 7A; species SO\sub{2}; reaction R781; the column-integrated change in
its rate, with thickness in centimetres, $8.2352\times10^{8}$, and in the old, wrong,
kilometre form, 8235.1671, kept to show the error (Section~\ref{sec:conv-col}).

\textbf{Example 4.} \texttt{e1.aflux.meanrel}: Experiment~1; actinic flux; mean relative
change. It is stored as the fraction $1.3553\times10^{-4}$ and printed as 0.013553\%. Keys
ending in \texttt{pct} hold percentages; keys ending in \texttt{rel} hold fractions.

% ===========================================================================
\section{The checker, and what it cannot catch}
\label{sec:checker}

The script \texttt{verify/check\_documents\_25m0328.py} finds every number in a document
written with three or more decimals, in scientific notation, or as a percentage with two
or more decimals, and accepts it only if a registry value, rounded to the number's own
significant digits, gives it (Section~\ref{sec:conv-numbers}). Three faults in the
checker (unbraced exponents, too tight a tolerance, DOIs read as numbers) were found
because it rejected correct numbers, and fixed. Chapters~1 to~10 of this report print
407 checked numbers and the twenty figures print 129 (checked by
\texttt{figures/check\_figures.py}); none is unmatched.

The checker proves that a number is in the registry, not that it came from the right
key. For example, \texttt{pnom.col\_dzw} holds $4.0775\times10^{8}$ per square
centimetre, which the checker would accept as the absorber column; but only
\texttt{pnom.column}, $4.0776\times10^{8}$, summed on the grid the absorber term uses,
reproduces $\tau_{\mathrm{req}}$. Every number in this report was therefore also
checked by hand against the key named beside it.

\textbf{Reproduction.} The eleven runs in the \texttt{conv} group, from the nominal run of
Experiment~2 to Experiment~11, all ended with \texttt{longdy} between 0.0097778 and
0.0099988, below the limit of 0.01. The registry can be rebuilt by anyone from the
committed output with \texttt{verify/recompute\_25m0328.py}.

% ===========================================================================
\section{One number, traced from end to end}
\label{sec:trace}

Take the 7.34 of Experiment~7, the number this report quotes most often, and follow it
from the model to the page.

\begin{enumerate}[leftmargin=*, itemsep=2pt]
  \item \textbf{The model output.} Three runs saved in the repository's \texttt{output}
    folder: the upward move, \texttt{Experiment7A\_UUV-upward18-25m0328.vul}; the
    downward move, \texttt{Experiment7B\_UUV-downward18-25m0328.vul}; and the reference,
    \texttt{Experiment2\_A1\_UUV-nominal-25m0328.vul}, with the absorber in its usual
    place. Each file holds the actinic flux field of its run.
  \item \textbf{The calculation.} \texttt{verify/recompute\_25m0328.py} loads the three
    files and calls its function \texttt{radiative()} twice, once for each move against
    the reference. The function subtracts the two fields and evaluates
    Equation~\eqref{eq:field}.
  \item \textbf{The registry.} The results are stored as \texttt{e7a.L2pct} and
    \texttt{e7b.L2pct}, and their ratio, 7.3367, as \texttt{e7.ratio\_7A\_over\_7B}.
  \item \textbf{The sentence.} Section~\ref{sec:e7} prints 2.2498\% and 0.3067\% with
    those keys, and the ratio as 7.34.
  \item \textbf{The checker.} It finds 2.2498\% and 0.3067\% in the chapter and accepts
    each, because the registry value rounded to the same digits gives it. The ratio,
    printed with two decimals, falls outside the checker's rule, so it was checked by
    hand: 2.2498 divided by 0.30665 is 7.3367, which rounds to 7.34.
\end{enumerate}

Any number in this report can be followed the same way, backwards from the key printed
beside it.

% ===========================================================================
\section{What the checks catch}
\label{sec:errors}

Recomputing every number and checking every printed value guards against three kinds of
error, and each can be seen in the registry or the code. \textbf{Units}: an integral over
altitude taken with the thickness in kilometres is $10^{5}$ times too small; the registry
keeps both forms of one such integral, \texttt{e7.7a.SO2.R781.col} and
\texttt{e7.7a.SO2.R781.col\_km}, so the size of the error can be seen
(Section~\ref{sec:conv-col}). \textbf{The wrong key}: a number can be in the registry
and still answer a different question, which is why every sentence is also checked
against the key it names (Section~\ref{sec:checker}). \textbf{Reading below the floor}:
a change in abundance smaller than 2.5278\% is not a measurement, so it is reported only
as an upper bound (Section~\ref{sec:conv-floor}).
